\documentclass[11pt]{article}
\usepackage[margin=2.5cm]{geometry}
\usepackage{amsmath}
\usepackage{mathptmx}
\usepackage[T1]{fontenc}
\begin{document}

\section*{Times through virtual fonts}

Paragraph timesone defines the \textbf{offset voltage} as the small output that stays when the input is
zero, and it keeps going so the line breaks twice; the ff in offset and buffer is a ligature that T1 Times
builds from two f glyphs.

Paragraph timestwo carries math whose italic letters Times shifts: $V_{GS} = V_{t} + V_{OD}$ and
$H(j\omega) = P(s)/T(s)$, with a frequency $f$ and a time $t$, long enough to wrap once more.

Paragraph timesthree is plain Times text without ligatures or math, so it patches the way any paragraph
in a Times document does, and it wraps over two lines.

\end{document}
